% Generated draft: requires mathematical and source review.
\documentclass[11pt]{article}
\usepackage[margin=25mm]{geometry}
\usepackage{fontspec}
\setmainfont{Latin Modern Roman}
\usepackage{amsmath,amssymb,mathtools,unicode-math}
\setmathfont{Latin Modern Math}
\usepackage{enumitem,needspace,xurl,hyperref}
\setlength{\parindent}{0pt}
\setlength{\parskip}{5pt}
\setlength{\emergencystretch}{3em}
\newcommand{\N}{\mathbb{N}}
\newcommand{\Z}{\mathbb{Z}}
\newcommand{\Q}{\mathbb{Q}}
\newcommand{\R}{\mathbb{R}}
\newcommand{\C}{\mathbb{C}}
\newcommand{\F}{\mathbb{F}}
\newcommand{\A}{\mathbb{A}}
\newcommand{\PP}{\mathbb{P}}
\newcommand{\E}{\mathbb{E}}
\newcommand{\eps}{\varepsilon}
\newcommand{\dd}{\,\mathrm{d}}
\begin{document}
\section*{UnsolvedMath ID: 30006444}
\textbf{Source problem code:} catalogue entry
\subsection{Definitions and mathematical statement}
For notation, if \(A\) is a commutative coefficient ring, a Laurent polynomial in two variables is a finite sum
\[
 f(x,y)=\sum_{(a,b)\in S}c_{a,b}x^a y^b,
 \qquad S\subset\Z^2\text{ finite},\quad c_{a,b}\in A.
\]
The supplied record does not specify the intended coefficient ring; this notation does not choose one for the conjecture. A monomial has a single nonzero term, with any further restriction on its coefficient dependent on the intended convention.

The source describes a Laurent polynomial as \emph{zero-mutable} if it permits a finite sequence of mutations ending in a monomial. Schematically, this means that there exist a finite integer \(r\geq 0\) and permitted mutations \(\mu_1,\ldots,\mu_r\) such that
\[
 f_0=f,\qquad f_j=\mu_j(f_{j-1})\quad(1\leq j\leq r),
 \qquad f_r\text{ is a monomial}.
\]
The permitted mutation operation, its admissibility conditions, and its coefficient conventions are not supplied. The displayed convention allows an empty sequence; the record does not explicitly discuss this boundary case.

In cluster algebra terminology, a \emph{cluster variable} is a mutable variable occurring in a seed obtained from an initial seed by a finite sequence of seed mutations. Interpreting a Laurent polynomial as such a variable also requires specifying the cluster algebra's coefficients and an identification with the relevant Laurent expressions.

\textbf{Conjecture as supplied, attributed to Gr\"afnitz--R.} Every zero-mutable Laurent polynomial in two variables is a cluster variable for a suitable quiver. Its quantifier structure is
\[
 \text{for every zero-mutable }f\text{ in two variables, there exists a suitable quiver }Q
 \text{ for which }f\text{ is a cluster variable}.
\]
The quiver may depend on \(f\). This is a formulation schema requiring clarification: the source record does not determine the coefficient domain, the Laurent-polynomial mutation operation, the allowed quivers and seed data, or the required identification of \(f\) with a cluster variable. In particular, two Laurent variables do not, from the supplied wording alone, impose a rank-two restriction on \(Q\). No normalization, positivity, support, or nondegeneracy hypothesis can be added from this record.

The original excerpt calls this a conjecture and adds ``maybe soon a theorem.'' The supplied dated curation labels it open; neither qualification constitutes independent verification here.
\subsection{Short English statement}
The proposed assertion is that every two-variable Laurent polynomial that can be mutated into a monomial occurs as a cluster variable for some suitable quiver. The coefficient conventions, allowed mutations, and meaning of the required cluster realization remain unspecified, so this is a statement-only draft requiring human mathematical review.
\subsection{Sources}
\begin{itemize}
\item Supplied local UnsolvedMath registry record, permanent ID \(30006444\), title \emph{Zero-Mutable Laurent Polynomials as Cluster Variables}, source problem code \texttt{catalogue entry}, in the Oberwolfach Reports Open Problems collection. The record gives proposed year \(2025\), leaves the proposer field empty, and attributes the conjecture within the original excerpt to ``Gr\"afnitz--R.'' The abbreviated attribution is retained without expansion.
\item Supplied source citation: \emph{Singularities (2026)}, Oberwolfach Reports, DOI \texttt{10.4171/owr/2025/45}, pp.~2413--2490. \url{https://doi.org/10.4171/owr/2025/45}. The citation's title year differs from the report year in the DOI and other supplied metadata; this discrepancy is unresolved.
\item Supplied literature source: \emph{Oberwolfach Report 45/2025}. \url{https://oa.tib.eu/renate/bitstreams/e0847358-34cd-4809-8330-e99d50f1e37e/download}.
\item Supplied literature source: \emph{Smoothings from zero-mutable Laurent polynomials}. \url{https://arxiv.org/abs/2503.18661}.
\item Literature assessment dated \texttt{2026-08-22}: the supplied curation labels the universal implication open and still conjectural, and describes related work as developing zero-mutable polynomials and proving important families of associated smoothing results, rather than the universal cluster-variable assertion. This is supplied dated curation, not an independently verified literature assessment. No linked source was visited for this draft.
\end{itemize}
\end{document}
